% !TeX root = ../main.tex

% 中英文摘要和关键字

\begin{abstract}
随着移动机器人在智能巡检、无人配送、灾害救援等复杂场景中的广泛应用，
对自主定位与环境感知能力提出了更高要求。传统激光或视觉同步定位与地图构建（SLAM）技术在
光照变化、环境退化、动态遮挡及重复场景等复杂条件下，容易出现定位精度下降、回环检测失败
及累积误差增大的问题。为提高机器人在复杂环境中的定位鲁棒性与环境感知能力，本文围绕多模
态信息辅助的机器人位姿优化开展研究，引入声源定位信息作为新的环境观测，为机器人定位提供
补充约束。

首先，针对复杂声学环境下声源定位精度不足的问题，研究了一种基于通道—频谱—时间特征建模的
声源定位方法。结合卷积神经网络、通道-频谱-时间特征重标定模块以及时序建模网络，对多通道
声学信号进行联合建模，提高了声源方向估计的准确性和鲁棒性，并通过公开数据集验证了所提方
法的有效性。其次，围绕机器人位姿优化问题，研究声源定位信息与激光惯性里程计系统的融合方
法，分析声学信息在回环检测、位姿约束及后端优化中的应用方式，探索声学观测与几何观测协同
优化机制，为复杂环境下机器人定位提供新的信息来源。最后，基于ROS平台构建多模态机器人定
位实验系统，在仿真环境及真实场景下开展实验，对声源定位性能、多模态融合定位精度及系统鲁
棒性进行综合验证。

本文研究工作实现了声学感知与机器人定位技术的有机结合，为多模态SLAM系统提供了一种新的
环境感知思路。研究结果表明，利用声源定位信息能够为机器人位姿估计提供有效补充，在复杂
环境下具有提升系统鲁棒性和定位精度的潜力，对智能机器人自主导航与多传感器融合感知具有
一定的理论意义和应用价值。
    \hdusetup{
        keywords = {关键词1、 关键词2、 关键词3、 关键词4、 关键词5},
    }
\end{abstract}

\begin{abstract*}
With the rapid development of intelligent inspection, autonomous delivery, and 
emergency rescue, mobile robots have placed higher demands on autonomous 
localization and environmental perception. Conventional simultaneous localization 
and mapping (SLAM) systems based on LiDAR or vision often suffer from degraded 
localization accuracy, unreliable loop closure detection, and accumulated drift 
in challenging environments with illumination changes, dynamic obstacles, 
repetitive structures, and feature degradation. To improve the robustness and 
accuracy of robot localization, this thesis investigates multi-modal 
information-assisted robot pose optimization by introducing sound source 
localization (SSL) as an additional environmental observation to complement 
conventional geometric measurements.

First, to address the limited robustness of sound source localization in complex 
acoustic environments, a sound source localization method based on channel-spectro-temporal 
feature modeling is studied. By integrating convolutional neural networks, channel-spectro-temporal 
feature recalibration, and temporal modeling networks, the proposed method jointly exploits 
multi-channel acoustic signals to improve the accuracy and robustness of direction-of-arrival 
estimation. The effectiveness of the proposed approach is validated on public benchmark datasets. 
Subsequently, a multi-modal fusion framework for robot pose optimization is investigated by 
incorporating acoustic observations into a LiDAR-inertial odometry system. The potential roles of 
acoustic information in loop closure detection, pose constraints, and back-end optimization are 
analyzed, and a collaborative optimization mechanism combining acoustic and geometric observations 
is explored to provide complementary information for robot localization in challenging environments. 
Finally, a multi-modal robot localization system is developed based on the Robot Operating System 
(ROS), and comprehensive experiments in both simulation and real-world environments are conducted 
to evaluate the performance of sound source localization, localization accuracy, and overall system 
robustness.

This work combines acoustic perception with robot localization and provides a new perspective 
for multi-modal SLAM systems. The results demonstrate that sound source localization has the 
potential to provide effective complementary constraints for robot pose estimation, thereby 
improving localization robustness and accuracy in complex environments. The proposed framework 
offers both theoretical significance and practical value for autonomous navigation and 
multi-sensor perception in intelligent robotic systems.


    \hdusetup{
        keywords* = {Robot localization, Simultaneous Localization and Mapping (SLAM), Sound Source Localization, Multi-modal Fusion, Pose Optimization},
    }
\end{abstract*}
